\chapter{Model Development and Results}

\section{Research Progression on IndustReal}

IndustReal was used as the controlled research setting for comparing temporal
heads, backbones, decoding, and feature fusion. The progression below records
the measured recognition results. Acc is frame accuracy, Edit is edit score,
and F1@50 is segmental F1 at an overlap threshold of 50 percent.

\begin{table}[H]
\centering
\scriptsize
\begin{tabularx}{\textwidth}{@{}Xrrr@{}}
\toprule
\textbf{Configuration} & \textbf{Acc} & \textbf{Edit} & \textbf{F1@50} \\
\midrule
v1 Huge-K710 + ASFormer (baseline) & 73.3 & 68.8 & 56.3 \\
v1 + Viterbi & 73.4 & 74.2 & 62.2 \\
v2 SSv2-giant + ASFormer + Viterbi & 74.0 & 77.3 & 66.5 \\
v2 SSv2-giant + DiffAct & 74.4 & 79.6 & 68.9 \\
Fusion (giant + IV2-B14) + ASFormer + Viterbi & 75.7 & 77.6 & 68.2 \\
\textbf{v4 Fusion + DiffAct --- best} & \textbf{84.21} & \textbf{87.19} & \textbf{79.50} \\
v4 + SAM3 enrichment & 74.9 & 79.5 & 70.1 \\
\bottomrule
\end{tabularx}
\caption{IndustReal research progression.}
\label{tab:industreal-results}
\end{table}

\section{Iteration-by-Iteration Interpretation}

The table is a development sequence rather than a list of unrelated
experiments. Each iteration answered a specific question about the system.

\begin{enumerate}
    \item \textbf{v1: VideoMAEv2-Huge with ASFormer.} This established the
    baseline for frame recognition and temporal segmentation. Its accuracy was
    usable, but the lower Edit and F1@50 scores showed that the predicted
    segments were not yet sufficiently aligned or stable at procedure
    boundaries.
    \item \textbf{v1 with Viterbi decoding.} The visual representation was
    unchanged; only the sequence decoding was improved. Edit increased from
    68.8 to 74.2 and F1@50 from 56.3 to 62.2. This demonstrated that temporal
    consistency was an important source of improvement, not just a larger
    backbone.
    \item \textbf{v2: SSv2-giant features.} Moving to the stronger giant
    representation raised accuracy and improved the overall feature quality.
    With ASFormer and Viterbi, the system reached 74.0 accuracy, 77.3 Edit,
    and 66.5 F1@50.
    \item \textbf{v2 with DiffAct.} Replacing the temporal head produced a
    further improvement in segment quality: F1@50 increased to 68.9 and Edit
    to 79.6. This identified DiffAct as the strongest offline temporal head in
    the tested progression.
    \item \textbf{Fusion of motion and appearance.} VideoMAEv2-giant and
    InternVideo2-B14 contributed different evidence. The fused representation
    with ASFormer and Viterbi reached 75.7 accuracy and 77.6 Edit, although its
    F1@50 did not exceed the v2 DiffAct result. This showed that feature fusion
    and temporal modeling had to be selected together.
    \item \textbf{v4: Fusion with DiffAct.} Combining the best representation
    with the strongest tested offline temporal head produced the clear best
    result: 84.21 accuracy, 87.19 Edit, and 79.50 F1@50. This became the
    reference architecture and feature recipe before the project moved to its
    own engine-model recordings.
    \item \textbf{SAM3 enrichment.} Adding segmentation-based enrichment did
    not improve the recognition result; the measured configuration fell to
    74.9 accuracy, 79.5 Edit, and 70.1 F1@50. SAM3 therefore remained an
    exploratory object/mask branch rather than a replacement for the temporal
    recognition path.
\end{enumerate}

\section{What the IndustReal Results Established}

The experiments established three practical conclusions. First, motion and
appearance information were complementary: motion described how the operator
was acting, while appearance helped distinguish visually similar parts and
contexts. Second, temporal decoding materially affected the quality of
procedure boundaries; a frame-accurate classifier was not sufficient by
itself. Third, the strongest offline result was not automatically the correct
deployment model. DiffAct benefited from sequence context that is unavailable
at the moment a live decision is required.

The research also exposed the limits of using a recognition model as a
correctness inspector. In the reported ablation, VideoMAE-only configurations
had 0\% incorrect-install recall, while the giant plus InternVideo2-B14 fusion
produced the first non-zero recall at 10.1\%. That improvement was useful
evidence that appearance fusion could help identify abnormal actions, but it
was not sufficient to claim reliable visual verification. The live
demonstrator therefore focuses on procedure-step recognition and flags
sequence deviations for review rather than certifying installation quality.

\section{Decision Before Moving to the Own Dataset}

\begin{table}[H]
\centering
\small
\begin{tabularx}{\textwidth}{@{}p{0.34\textwidth}X@{}}
\toprule
\textbf{Research decision} & \textbf{How it shaped the next phase} \\
\midrule
Retain dual-stream features & The engine-model pipeline used the complementary
motion and appearance encoders rather than reverting to a single backbone. \\
Use v4 as the offline reference & The fused representation plus DiffAct
provided the upper-bound reference for the in-house dataset. \\
Do not deploy DiffAct unchanged & Its offline context requirement led to a
causal temporal head for live inference. \\
Keep decoding explicit & Fixed-lag sequence commitment was retained so that
the interface would report stable steps rather than every transient class
change. \\
Treat correctness as a separate problem & The first deployment target remained
step recognition, out-of-order detection, and procedure-state presentation. \\
\bottomrule
\end{tabularx}
\caption{Research decisions carried into the engine-model development phase.}
\label{tab:research-decisions}
\end{table}

\section{Evaluation Metrics and Selection Logic}

Accuracy measures the proportion of correctly classified frame positions. Edit
score measures the similarity of the predicted and reference action order after
collapsing repeated labels, so it is sensitive to extra, missing, or misplaced
segments. F1@50 measures whether predicted action segments overlap reference
segments by at least 50\%. The progression was selected using all three: a
model was not considered better merely because it improved frame accuracy if
its procedure boundaries became less reliable.

\section{Engine-Model Dataset Results}

\begin{table}[H]
\centering
\small
\begin{tabularx}{\textwidth}{@{}Xrrr@{}}
\toprule
\textbf{Configuration} & \textbf{Acc} & \textbf{Edit} & \textbf{F1@50} \\
\midrule
v4 Fusion + DiffAct (offline) & 96.21 & 97.33 & 97.51 \\
Baseline (online) & 92.96 & 75.65 & 81.36 \\
\textbf{Final Online V1} & \textbf{94.12} & \textbf{84.45} & \textbf{89.50} \\
\bottomrule
\end{tabularx}
\caption{Engine-model dataset results supplied for the project evaluation.}
\label{tab:engine-results}
\end{table}

The offline result is a useful upper-bound reference: it benefits from
sequence context that is unavailable at deployment time. Final Online V1
improves substantially over the online baseline in temporal quality while
retaining strong frame accuracy. The gap between offline and online results is
expected and is the reason the deployment metrics are reported separately.

\section{From Offline Models to Causal Inference}

The causal temporal convolutional network uses nine dilated blocks with 128
feature maps and approximately 875 thousand trainable parameters. Its
receptive field is approximately 205 seconds, but its computation is causal:
the output at time $t$ depends on the current and previous feature sequence,
not on future frames. This gives the model enough history to understand
procedure context while preserving the online constraint.

The final class set contains eleven classes: the ten modeled parts and a
background state. The network generates class scores for each feature update.
The fixed-lag decoder then converts those scores into stable committed steps for
the application.

\section{Robustness Iteration and Acceptance Gate}

The deployed checkpoint, \texttt{robust\_v1}, was trained with normal recordings,
mid-start crops, and scrambled-order variants. The project acceptance gate
measured both normal behavior and the harder scrambled-order behavior.

\begin{table}[H]
\centering
\small
\begin{tabularx}{\textwidth}{@{}p{0.34\textwidth}rrrX@{}}
\toprule
\textbf{Checkpoint / run} & \textbf{Normal} & \textbf{Scrambled} &
\textbf{Cold start} & \textbf{Decision} \\
\midrule
Previous baseline & 93.45\% & 72.58\% & 19.2 s & Reference only \\
Fine-tune A & 92.39\% & 91.40\% & --- & Rejected: normal floor \\
Fine-tune A2 & 92.95\% & 88.84\% & --- & Passed interim check; not selected \\
\texttt{robust\_v1} & 93.06\% & 90.57\% & 0.2 s & Deployed checkpoint \\
\bottomrule
\end{tabularx}
\caption{Robustness iteration and deployment selection.}
\label{tab:robustness-gate}
\end{table}

\section{Checkpoint-to-Application Handoff}

Once the checkpoint passed the gate, the model was not treated as an isolated
file. The feature-extraction path, class vocabulary, decoder lag, and service
state all had to remain aligned. The replay service loads the checkpoint and
runs it over a complete recording; the live service loads the same checkpoint
and updates it incrementally as new clips arrive. The shared state contract
then allows the browser and HoloLens clients to present the result without
duplicating model logic.

\begin{table}[H]
\centering
\small
\begin{tabularx}{\textwidth}{@{}p{0.33\textwidth}X@{}}
\toprule
\textbf{Handoff step} & \textbf{Engineering purpose} \\
\midrule
Prepare or compute fused features & Guarantee the same clip length, sampling
rate, normalization, and feature ordering used during training. \\
Load the causal checkpoint & Use the selected temporal head and class order
consistently in replay and live services. \\
Run fixed-lag commitment & Convert repeated model scores into stable step
transitions and measurable delay. \\
Publish procedure state & Give presentation clients current step, progress,
and attention information through one interface. \\
\bottomrule
\end{tabularx}
\caption{The engineering handoff from a trained checkpoint to the
demonstration.}
\label{tab:checkpoint-handoff}
\end{table}

\section{Optional Visual Enrichment}

SAM3-based enrichment was explored as a supporting visual branch for masks and
part inventory. It was useful for investigating object-level cues, but it did
not replace the temporal recognition path and did not become the basis of the
deployed live checkpoint. The separation keeps the demonstrated system focused:
the primary model recognizes procedure state from the egocentric sequence.
